\section{WARUNKI ZWYCIĘSTWA}

Ogólna zasada:

Istnieją dwa rodzaje zwycięstwa: zwycięstwo zespołowe i zycięstwo indywidualne.
Grupa, tzn. oddział wychodzi z gry zwycięsko jeżeli odnajdzie Czarne Wrota, zniszczy je i co najmniej jedna postać wyjdzie żywa z labiryntu.
Może to wymagać rozegrania więcej niż jednej przygody tzn. więcej niż jednej wizyty w labiryncie.

Zwycięzcą indywidualnym zostaje ta postać, która ujdzie z życiem ze wszystkich rozegranych przygód, zdobędzie najwięcej punktów doświadczenia oraz zgromadzi największy majątek w sztukach złota i kosztownościach.
Jeżeli gracze decydują się na zakończenie gry przed znalezieniem lub zniszczeniem Czarnych Wrót zwycięstwo grupowe jest niemożliwe

Sposób postępowania:

Oddział wchodzi do labiryntu i poszukuje Czarnych Wrót.
Jeżeli w jakimś momencie gracze decydują, że istnieje zbyt duże prawdopodobieństwo, że oddział zostanie wybity do nogi zanim je odnajdzie, może on wyjść z labiryntu tą samą drogą, którą doń wszedł, tzn. przez bramę labiryntu na pierwszym poziomie.
Taki fragment gry – od chwili wejścia do labiryntu do chwili opuszczenia go bez wykonania głównego zadania – nazywany jest przygodą.
Po wyleczeniu wszystkich ran oddział może ponownie wejść do labiryntu.
Zbadane podczas poprzednich przygód segmenty są znane w przygodach następnych.

Zasady szczegółowe:

[17.1] Sztuki złota i kosztowności, zdobyte przez postać w trakcie jednej przygody są zostawiane w domu przed rozpoczęciem następnej, o ile nie zostały wydane na zwiększenie poszczególnych współczynników (patrz 12.2).

[17.2] Potwory, opanowane przy pomocy czaru ,,Urok'' zostają wyzwolone spod jego władzy w momencie wyjścia oddziału z labiryntu.

[17.3] Po każdej przygodzie, w trakcie której zostanie zabity jeden (lub więcej) członek oddziału można wprowadzić do gry nowe postacie, tak, by oddział liczył 3 bohaterów i 3 uczniów.
Nowo wprowadzone postacie zaczynają oczywiście z zerowym kontem punktów doświadczenia i skarbów.

[17.4] Przed każdym kolejnym wejściem do labiryntu postać może wybrać nową broń, jednak tak, by miała przy sobie nie więcej niż dwa jej rodzaje.

[17.5] Zakłada się, że między wyjściem z labiryntu a ponownym do niego wejściem upływa 6 miesięcy.
W ciągu tego czasu postacie nie tylko leczą wszystkie odniesione rany, ale także ,,starzeją się'', a więc wolniej nabywają nowe, wyższe umiejętności.
Jest to odzwierciedlone w rosnącym koszcie nabywania punktów, zwiększających poszczególne współczynniki.
Koszty te rosną o 10% po zakończeniu każdej kolejnej przygody, poczynając od drugiej (tzn. od drugiego wyjścia z labiryntu).
